The constraints can be written as $Ax=b$ with
\[
  A=\begin{pmatrix}
  1 & 1 & 1 & 1 \\
  1 & 0 & -1 & 0 \\
  2 & 1 & 0 & 1 
\end{pmatrix}, \qquad b = \begin{pmatrix}
1\\ -1\\ 0
  \end{pmatrix}.
  \]
  The system is compatible as
  \[
x = \cvect{0 \\ 0 \\ 1 \\ 0}
  \]
  is feasible.

In principle, in the presence of three constraints, we should select 3
variables out of the 4 to be basic variables. But for each of the 4
possibilities, the corresponding basic matrix is not invertible,
precluding from eliminating the constraints. The reason is that the
constraint matrix is not of full rank.  The third constraint happens
to be the sum of the two other constraints. It is therefore redundant
and must be removed. We obtain the system with the following data:
\[
A=\begin{pmatrix}
  1 & 1 & 1 & 1 \\
  1 & 0 & -1 & 0
\end{pmatrix}, \qquad b = \begin{pmatrix}
1\\ -1
\end{pmatrix}.
\]
As we have two constraints, we can eliminate two variables. There are
6 ways to do it. But because $x_1$ cannot be in the basis, we need to
select two variables out of three.  And there are three ways to do it.

\begin{enumerate}
\item $x_2$ and $x_3$ in the basis. The basic matrix is obtained  by extracting columns $2$
and $3$ from $A$:
\[
B= \begin{pmatrix}
  1 & 1 \\
  0 & -1
\end{pmatrix}
\]
which is  invertible and happens to be equal to its own inverse:
\[
B^{-1}= \begin{pmatrix}
  1 & 1 \\
  0 & -1
\end{pmatrix}.
\]
In order to apply the generic formula
 $$AP= \begin{pmatrix}
B & N
\end{pmatrix},$$
we need to permute the
columns of the matrix so that columns 2 and 3 become the two first
columns of the permuted matrix. We use the following permutation
matrix, 
\[
P = \begin{pmatrix}
0 & 0 & 1 & 0  \\
1 & 0 & 0 & 0  \\
0 & 1 & 0 & 0  \\
0 & 0 & 0 & 1  \\
\end{pmatrix},
\]
so that $PP^T=I$ and 
\[
AP = \begin{pmatrix}
  1 & 1 & 1 & 1  \\
  0 & -1 & 1 & 0 
\end{pmatrix}.
\]
The constraints can now be written
as
\[
Ax = AP(P^Tx) = Bx_B + Nx_N = b,
\]
and we obtain
\begin{align*}
x_B = \begin{pmatrix}
x_2\\ x_3
\end{pmatrix} & = B^{-1}(b-Nx_n)\\
& = \begin{pmatrix}
1 & 1 \\ 0 & -1
\end{pmatrix}
\left(
\rvect{1 \\ -1} -
\begin{pmatrix}
  1 & 1 \\ 1 & 0
\end{pmatrix}
\cvect{x_1\\x_4}
\right) \\
& = \cvect{-2 x_1 - x_4 \\ x_1 + 1}.
\end{align*}
The basic variables are $x_2$ and $x_3$, and the non basic variables
are $x_1$ and $x_4$.
\item $x_2$ and $x_4$ in the basis. The basic matrix is obtained by
  extracting columns $2$ and $4$ from $A$:
\[
B= \begin{pmatrix}
  1 & 1 \\
  0 & 0
\end{pmatrix}.
\]
This matrix is not invertible, and the elimination of constraints
cannot be performed. 
\item $x_3$ and $x_4$ in the basis. The basic matrix is obtained by
  extracting columns $2$ and $4$ from $A$:
\[
B= \begin{pmatrix}
  1 & 1 \\
  -1 & 0
\end{pmatrix}.
\]
It is invertible, and 
\[
B^{-1}= \begin{pmatrix}
  0 & -1 \\
  1 & 1
\end{pmatrix}.
\]
In order to apply the generic formula
 $$AP= \begin{pmatrix}
B & N
\end{pmatrix},$$
we need to permute the
columns of the matrix so that columns 3 and 4 become the two first
columns of the permuted matrix. We use the following permutation
matrix, 
\[
P = \begin{pmatrix}
0 & 0 & 1 & 0  \\
0 & 0 & 0 & 1  \\
1 & 0 & 0 & 0  \\
0 & 1 & 0 & 0  \\
\end{pmatrix},
\]
so that $PP^T=I$ and 
\[
AP = \begin{pmatrix}
  1 & 1 & 1 & 1  \\
  -1 & 0 & 1 & 0 
\end{pmatrix}.
\]
The constraints can now be written
as
\[
Ax = AP(P^Tx) = Bx_B + Nx_N = b,
\]
and we obtain
\begin{align*}
x_B = \begin{pmatrix}
x_3\\ x_4
\end{pmatrix} & = B^{-1}(b-Nx_n)\\
& = \begin{pmatrix}
0 & -1 \\ 1 & 1
\end{pmatrix}
\left(
\rvect{1 \\ -1} -
\begin{pmatrix}
  1 & 1 \\ 1 & 0
\end{pmatrix}
\cvect{x_1\\x_2}
\right) \\
& = \cvect{1 + x_1 \\-2 x_1 - x_2 }.
\end{align*}
The basic variables are $x_3$ and $x_4$, and the non basic variables
are $x_1$ and $x_2$.

\end{enumerate}
